\chapter{Другий вихід: як мозок керує хімією тіла}
\chaptermark{Хімічний вихід}
\label{sec:chemoutput}
\begin{chaptergoals}
\item як газ \nt{NORA} і гальмо \nt{ACET} підлаштовують ритмоводій і чому гальмо швидше;
\item як кров працює широкомовною шиною, а \nt{ADRE} переводить тіло в режим бігу;
\item чому каскад гормонів зі зворотним зв'язком не може бути водночас точним і стійким: $K<8$;
\item як утомлений рецептор читає частоту порцій і як добовий ритм підстроюється за світлом.
\end{chaptergoals}

\section{Швидкий і повільний виходи}

У розділі~\ref{sec:muscles} ми бачили швидкий вихід машини --- м'язи. Але є в неї й другий, повільний. Мозок тримає в потрібних межах температуру тіла, частоту серця, кількість води та цукру, готує тіло до бігу чи до сну. Для цього він не рухає суглобами, а випускає хімічні сигнали. Шляхів два. Перший --- нерви, що йдуть просто до внутрішніх органів: до серця, судин, кишківника, залоз. Другий --- кров: деякі нейрони й залози викидають речовину не в щілину до сусіда, а просто в кровотік.

Кров --- найширокомовніша шина з усіх, що траплялися в книзі. Широкомовні канали розділів~\ref{sec:serocomputer}--\ref{sec:acet} розсилали сигнал на великі ділянки мозку; кров розсилає його кожній клітині тіла. Кров обходить тіло за час порядку хвилини, тож повідомлення доходить до всіх за десятки секунд і тримається хвилини й години, доки речовина не зруйнується. Адреси в такої розсилки немає зовсім. Як і в твердженні~\ref{prop:receptor}, зміст повідомлення обирає отримувач: відгукуються лише клітини, що мають рецептор до цієї речовини.

\section{Газ і гальмо: два проводи до органів}

Найкращий приклад --- серце. У серці є власний ритмоводій, як у \nt{SERO}-нейрона розділу~\ref{sec:serocomputer}: він сам задає ритм, і якщо відключити всі нерви, серце б'ється приблизно сто разів на хвилину. Мозок ритму не породжує, а лише підлаштовує його, до того ж двома проводами протилежного знака (рис.~\ref{fig:gasbrake}). Одним іде \nt{NORA}: це \emph{газ}, він прискорює ритм. Другим --- \nt{ACET}: це \emph{гальмо}, він сповільнює. Грубо,
\begin{empheq}[box=\keybox]{equation}
f \;\approx\; f_0 \;+\; a_S\,S \;-\; a_P\,P ,
\label{eq:gasbrake}
\end{empheq}
де $f_0\approx100$ ударів на хвилину --- власний ритм ритмоводія, $S$ і $P$ --- активність газу й гальма. У спокої гальмо натиснуте, і серце б'ється зазвичай порядку шістдесяти--сімдесяти разів на хвилину; під навантаженням гальмо відпускають і натискають газ.

\begin{figure}[t]
\centering
\begin{tikzpicture}[>=Stealth,
  blk/.style={draw,very thick,rounded corners=2pt,align=center,font=\small},
  pace/.style={circle,draw,very thick,double,double distance=1.2pt,minimum size=1.8cm,inner sep=0pt,font=\scriptsize,fill=red!8,align=center},
  inh/.style={-{Bar[width=7pt]},thick,red!70!black}]
\node[blk,fill=blue!5,text width=1.6cm] (B) at (0,0) {команда\\мозку};
\node[pace] (H) at (6.2,0) {ритмо-\\водій};
\draw[->,very thick,orange!85!black] (B.north east) to[bend left=20] node[above,font=\scriptsize,black,align=center] {\nt{NORA}: газ, секунди} (H.north west);
\draw[inh,very thick,blue!60!black] (B.south east) to[bend right=20] node[below,font=\scriptsize,black,align=center] {\nt{ACET}: гальмо, частки секунди} (H.south west);
\draw[->,very thick] (H) -- (8.4,0) node[right,font=\small] {частота $f$};
\node[blk,fill=orange!10,text width=1.6cm] (A) at (0,-2.6) {\nt{ADRE}\\у кров};
\draw[->,thick,orange!85!black] (B) -- (A);
\foreach \y/\lab in {-2.0/{серце},-2.6/{судини},-3.2/{печінка, м'язи}}{
  \draw[->,thick,densely dashed,orange!85!black] (A.east) -- (4.3,\y) node[right,font=\scriptsize,black] {\lab};}
\end{tikzpicture}
\caption{Два проводи протилежного знака до ритмоводія серця: газ \nt{NORA} повільний, гальмо \nt{ACET} швидке. Унизу --- той самий газ широкомовною шиною: \nt{ADRE}-клітини викидають адреналін у кров, і сигнал отримують усі органи з відповідним рецептором (штрихові стрілки).}
\label{fig:gasbrake}
\end{figure}

Це двопровідний код розділу~\ref{sec:glutcomputer} і пара м'язів розділу~\ref{sec:muscles}, тільки тепер проводи керують не суглобом, а темпом ритмоводія. І швидкість у проводів різна. Обидва працюють як фільтри нижніх частот, але гальмо пропускає зміни в кілька разів швидше за газ~\cite{Berger1989}: у \nt{ACET} шлях короткий: пов'язаний із рецептором білок сам відкриває калієвий канал, без другого посередника всередині клітини~\cite{Logothetis1987gbg}, а \nt{NORA} діє через ланцюжок реакцій усередині клітини. Інженер упізнає тут прийом із двома приводами різної швидкості: швидкий робить точне миттєве підлаштування, повільний --- великі тривалі зміни.

Тут же знаходить свою роль \nt{ADRE}, про який у таблиці~\ref{tab:types} було написано, що його роль у мозку неясна. Поза мозком вона ясна. Частина клітин газового шляху не посилає провід до органа, а викидає адреналін просто в кров. Це той самий газ, але широкомовною шиною: за десятки секунд він доходить до серця, судин, печінки й м'язів одразу і тримається хвилини. Адресний газ \nt{NORA} підлаштовує один орган; широкомовний газ \nt{ADRE} переводить у режим бігу все тіло.

\section{Каскад зі зворотним зв'язком}

Багато гормонів випускає не одна залоза, а каскад. Невелика група нейронів в основі мозку викидає в кров перший сигнал; він змушує залозу-посередника випустити другий; той --- кінцеву залозу випустити гормон, який і діє на тіло. А кінцевий гормон повертається до перших ступенів і гальмує їх. Виходить підсилювач із трьох каскадів, охоплений від'ємним зворотним зв'язком, --- регулятор рівня, як у \nt{SERO}-системи у формулі~\eqref{eq:feedback}.

Опишімо кожен ступінь RC-колом зі сталою часу $\tau$ і підсиленням $g$, а зворотний зв'язок --- коефіцієнтом $k$:
\begin{equation}
\tau\,\frac{\dd x_1}{\dd t} = -x_1 + g\bigl[u - k\,x_3\bigr]_+ ,\qquad
\tau\,\frac{\dd x_2}{\dd t} = -x_2 + g\,x_1 ,\qquad
\tau\,\frac{\dd x_3}{\dd t} = -x_3 + g\,x_2 ,
\label{eq:cascade}
\end{equation}
де $u$ --- команда мозку, $x_3$ --- рівень кінцевого гормону. Підсилення петлі $K=g^3k$. У рівновазі $x_3=g^3u/(1+K)$, і, як у будь-якого регулятора зі зворотним зв'язком, велика петля робить рівень нечутливим до розкиду підсилень ступенів. Але кожен ступінь зсуває фазу: на частоті $\omega=\sqrt3/\tau$ три однакові ступені разом дають зсув на пів оберту й ослаблення у вісім разів. Звідси класичний результат теорії автоматичного керування:
\begin{empheq}[box=\keybox]{equation}
K \;<\; 8 ,
\label{eq:cascadestable}
\end{empheq}
інакше регулятор починає коливатися з періодом близько $2\pi\tau/\sqrt3\approx3.6\,\tau$.

\begin{keyblock}
\begin{proposition}[Точність проти стійкості]
\label{prop:cascade}
У триступеневого каскаду з пропорційним зворотним зв'язком похибка рівня дорівнює $1/(1+K)$, а стійкий він лише за $K<8$. Тому точніше, ніж приблизно на десять відсотків, такий регулятор рівня не тримає; спроба підняти підсилення перетворює його на генератор коливань.
\end{proposition}
\end{keyblock}

Ми перевірили це на моделі~\eqref{eq:cascade} (рис.~\ref{fig:cascade}). За $K=4$ рівень після ввімкнення трохи коливається й усталюється. За $K=7$ він теж усталюється, але повільно. За $K=12$ коливання не згасають і не ростуть без меж: ступінь не може видати від'ємний сигнал, і каскад виходить на рівні пульсації між $0.83$ і $1.24$ від середнього рівня з періодом $3.7$--$3.9\,\tau$.

\begin{figure}[t]
\centering
\begin{tikzpicture}[>=Stealth,x=0.3cm,y=1.2cm]
\draw[->,gray!70!black] (0,0) -- (41.5,0) node[right,font=\small,black] {$t/\tau$};
\draw[->,gray!70!black] (0,0) -- (0,2.8) node[left,font=\small,black] {$x_3/x_3^*$};
\foreach \t in {0,10,20,30,40}{\draw[gray!60!black] (\t,-0.05) -- (\t,0.05); \node[below,font=\scriptsize] at (\t,0) {\t};}
\foreach \f in {1,2}{\draw[gray!60!black] (-0.3,\f) -- (0.3,\f); \node[left,font=\scriptsize] at (0,\f) {\f};}
\draw[densely dashed,gray!60!black] (0,1) -- (40,1);
\draw[thick,blue!60!black] plot file {figures/cascade_K4.dat};
\draw[thick,red!70!black] plot file {figures/cascade_K12.dat};
\node[font=\scriptsize,blue!60!black,anchor=west] at (16,0.55) {$K=4$: усталюється};
\node[font=\scriptsize,red!70!black,anchor=west] at (16,1.75) {$K=12$: пульсує};
\end{tikzpicture}
\caption{Каскад із трьох ступенів зі зворотним зв'язком за~\eqref{eq:cascade}, рівень кінцевого гормону в частках рівноважного $x_3^*$. Коли підсилення петлі менше за вісім, рівень усталюється; коли більше --- каскад сам породжує рівні пульсації.}
\label{fig:cascade}
\end{figure}

Справжні гормони каскадів справді виходять пульсами: рівень гормону стресу, наприклад, коливається з періодом близько години. За однією з гіпотез, ці пульсації породжує сам зворотний зв'язок між ступенями із запізненням, і окремий генератор ритму для них не потрібен~\cite{Walker2010}. Це та сама причина коливань, що й у петлі \nt{GLUT}--\nt{GABA} у твердженні~\ref{prop:ei} і в петлі розтягу м'яза в умові~\eqref{eq:delaystable}: від'ємний зворотний зв'язок, який запізнюється.

\section{Імпульсний код гормонів}

Пульси бувають не лише побічним ефектом, а й кодом. Нейрони, що керують розмноженням, випускають свій сигнал у кров короткими порціями приблизно раз на годину. Якщо залозі-отримувачу подавати той самий сигнал не порціями, а безперервно, вона не працює на повну силу, як за пульсів, а замовкає; якщо знову подавати порціями раз на годину, робота відновлюється~\cite{Belchetz1978}. Отже, отримувач читає не рівень, а \emph{частоту} порцій.

Для інженера тут немає загадки. Рецептор, який за тривалої дії речовини втомлюється й перестає відповідати, --- фільтр верхніх частот, як синапс, що виснажується~\eqref{eq:depression}, з розділу~\ref{sec:code}. Безперервний сигнал він глушить, а зміни пропускає. Такому отримувачеві можна щось повідомити лише імпульсами, і інформація переходить із рівня в частоту й форму порцій. Різні частоти порцій до того ж по-різному діють на різні клітини однієї й тієї самої залози, тож однією хімічною лінією можна передати більше ніж одну величину --- як одним проводом \nt{DOPA} в розділі~\ref{sec:dofaalt} ішли швидка й повільна складові.

\section{Добовий ритм: повільний генератор режимів}

Найповільніший ритм тіла --- добовий. Його породжує від'ємний зворотний зв'язок із запізненням усередині клітин: гени виробляють білки, які із затримкою в кілька годин пригнічують власне виробництво~\cite{TakahashiJS2017}. Знову запізнілий від'ємний зворотний зв'язок --- вчетверте в цій книзі, --- і знову ритм, цього разу з періодом близько доби. Головний генератор --- невелика група з десятків тисяч нейронів в основі мозку, ритм яких синхронізує решту клітин тіла.

Власний період цього генератора в людини не рівно доба, а в середньому $24.18$ години~\cite{Czeisler1999}. Щодня його підводить світло, що надходить від датчиків ока (розділ~\ref{sec:sensors}). Для електронника це \emph{фазове автопідстроювання частоти}: генератор із власною частотою $\omega_0$ підлаштовується до зовнішнього сигналу з частотою $\omega$, і різниця фаз $\varphi$ підпорядковується рівнянню
\begin{empheq}[box=\keybox]{equation}
\frac{\dd\varphi}{\dd t} \;=\; (\omega-\omega_0) \;-\; K_\varphi\,\sin\varphi .
\label{eq:pll}
\end{empheq}
Підстроювання утримується, якщо $|\omega-\omega_0|<K_\varphi$. Генераторові з періодом $24.18$ години треба щодня зсуватися приблизно на одинадцять хвилин; ранкового світла для такого зсуву зазвичай вистачає. У полярну ніч або в печері без світла підстроювання немає, і ритм переходить на власний період.

Добовий генератор --- не тактовий генератор комп'ютера. Він не відлічує кроків обчислень і не синхронізує імпульсів нейронів. Він перемикає \emph{режими}: сон і неспання, рівні гормонів, температуру тіла, готовність до їжі. Це повільний широкомовний регістр того самого роду, що й регістри розділів~\ref{sec:serocomputer}--\ref{sec:acet}, тільки його значення змінюється за розкладом, підведеним до сонця.

\begin{keyblock}
\begin{proposition}[Хімічний вихід]
\label{prop:chemoutput}
Повільний вихід машини збудовано з тих самих деталей, що й швидкий. Органи отримують два проводи протилежного знака й різної швидкості --- газ \nt{NORA} і гальмо \nt{ACET}, --- а весь організм одразу --- широкомовні сигнали кров'ю, зокрема адреналін \nt{ADRE}. Рівні гормонів тримають каскади з від'ємним зворотним зв'язком, точність яких обмежена стійкістю; частина гормонів кодується частотою порцій; добовий ритм --- повільний генератор із фазовим підстроюванням за світлом. Будову шляхів і досліди з порціями та з періодом виміряно; пояснення пульсацій каскаду самим зворотним зв'язком --- гіпотеза.
\end{proposition}
\end{keyblock}

\section{Підсумок}

\begin{table}[t]
\centering
\footnotesize
\setlength{\tabcolsep}{4pt}
\renewcommand{\arraystretch}{1.15}
\begin{tabular}{>{\raggedright}p{3.1cm}>{\raggedright}p{4.8cm}>{\raggedright\arraybackslash}p{4.9cm}}
\toprule
Що робить хімічний вихід & Як & Що відомо \\
\midrule
підлаштовує органи & газ \nt{NORA} і гальмо \nt{ACET}~\eqref{eq:gasbrake} & виміряно; гальмо швидше за газ \\
переводить усе тіло в режим бігу & адреналін \nt{ADRE} у кров & виміряно \\
тримає рівні гормонів & каскад зі зворотним зв'язком, $K<8$~\eqref{eq:cascadestable} & каскади й зворотний зв'язок виміряно; пульсації від самої петлі --- гіпотеза \\
передає частотою порцій & утомлений рецептор як фільтр верхніх частот & залежність від порцій виміряно \\
перемикає режими за добу & генератор із фазовим підстроюванням за світлом~\eqref{eq:pll} & період і підстроювання світлом виміряно \\
\bottomrule
\end{tabular}
\caption{Як машина керує хімією тіла.}
\label{tab:chemoutput}
\end{table}

У машини два виходи (таблиця~\ref{tab:chemoutput}). Швидкий --- м'язи: мілісекунди, адресно, у кожен суглоб. Повільний --- хімія тіла: секунди, хвилини й доба, двома проводами до органів і кров'ю до всіх одразу. Деталі в обох виходів ті самі --- два проводи протилежного знака, ритмоводії, зворотний зв'язок та його запізнення, --- і ті самі, з яких зібрано машину всередині.

\section{Задачі}

\begin{exercise}[\lvlA]
\label{ex:16:1}
\emph{Кров як фільтр.} Гормон виводиться з крові з періодом напіврозпаду $t_{1/2}=10$~хв і випускається короткими порціями кожні $T=60$~хв. У скільки разів максимум концентрації більший за мінімум і у скільки разів ослаблена основна гармоніка порцій? За якого $t_{1/2}$ максимум перевищує мінімум лише вдвічі?
\end{exercise}

\begin{exercise}[\lvlA]
\label{ex:16:2}
\emph{Добове автопідстроювання.} Власний період генератора~\eqref{eq:pll} дорівнює $24.18$~год, світло зсуває фазу не більше ніж на годину за добу: $K_\varphi=2\pi/24$~рад/доб. Знайдіть усталену різницю фаз $\varphi^*$ у хвилинах і час релаксації. За скільки діб розузгодження після стрибка зовнішнього сигналу на $\pm6$~год стане меншим за годину?
\end{exercise}

\begin{exercise}[\lvlB]
\label{ex:16:3}
\emph{Точність проти стійкості.} Лінеаризований каскад~\eqref{eq:cascade} подовжили до $n$ однакових ступенів. Знайдіть критичне підсилення $K_c(n)$ і найменшу досяжну похибку $1/(1+K_c)$ за $n=3$ і $n=6$. До чого прямує ця похибка за $n\to\infty$?
\end{exercise}

\begin{exercise}[\lvlB]
\label{ex:16:4}
\emph{Газ і гальмо.} У~\eqref{eq:gasbrake} внески газу й гальма --- виходи RC-кіл із $\tau_S=2$~с і $\tau_P=0.4$~с, команди не більші за $80$ і $60$ ударів на хвилину. Ритм $f_0=100$; у спокої гальмо $35$, газ $0$, $f=65$. За який час можна вийти на $f=120$? А якщо гальма не відпускати й працювати лише газом?
\end{exercise}

\begin{exercise}[\lvlB]
\label{ex:16:5}
\emph{Моделювання: каскад із затримкою.} У зворотний зв'язок функції \texttt{cascade} з \texttt{models/\allowbreak chem\_models.py} (скопіюйте її, файл не запускайте) додайте затримку $d$: перший ступінь бачить $x_3(t-d)$. Знайдіть $K_c$ і період на межі за $d/\tau=0$, $0.5$, $1$, $2$ і перевірте моделюванням за $0.9K_c$ і $1.1K_c$.
\end{exercise}

\begin{exercise}[\lvlB]
\label{ex:16:6}
\emph{Отримувач, що читає порції.} Рецептор утомлюється: $y=[c-a]_+$, $\tau_a\,\dd a/\dd t=c-a$, $\tau_a=30$~хв. Гормон надходить порціями по $6$~хв із періодом $T$ за тієї самої середньої дози $\bar c$. Знайдіть середній відгук за $T=15$, $60$, $120$~хв і $T\to\infty$. Який він за сталої концентрації $\bar c$?
\end{exercise}

\begin{exercise}[\lvlC]
\label{ex:16:7}
\emph{Генератор чи зворотний зв'язок?} Запропонуйте дослід, який відрізнить пульсації, породжені запізнілим зворотним зв'язком каскаду~\eqref{eq:cascade}, від окремого ритмоводія. Чи можна розрізнити гіпотези за зсувом періоду, якщо $K$ вдається змінити лише в півтора раза, а період вимірюється з точністю $5\%$?
\end{exercise}
